\subsection{Perplexity 与长度偏差：两个看似简单却容易误读的指标}

slides 75--77 是问答后的评估补充。perplexity 把平均 token 交叉熵指数化，直觉上表示模型在每个位置“等效犹豫多少个候选”；它仍是预训练开发的重要监控指标，但不直接衡量 instruction following、事实性或 Agent 成功。随后课程回到 LLM judge 的 spurious correlation，以输出长度偏好说明评估器会奖励与真实目标无关的特征。

\lecturefigure{slides-images/slide-075.png}{官方 slide 75：perplexity 的公式与等效候选数直觉}
\lecturefigure{slides-images/slide-076.png}{官方 slide 76：语言模型 perplexity 的历史下降与开发用途}
\lecturefigure{slides-images/slide-077.png}{官方 slide 77：LLM judge 长度偏差与 length-controlled correction}

\begin{importantbox}{Perplexity 定义与符号}
若长度为 $L$ 的 token 序列为 $x_{1:L}$，平均自然对数交叉熵为
\[
H=-\frac{1}{L}\sum_{i=1}^{L}\log p(x_i\mid x_{<i}),
\qquad \mathrm{PPL}=\exp(H).
\]
其中 $x_i$ 是第 $i$ 个真实 token，$x_{<i}$ 是前文，$p$ 是模型给真实 token 的概率。若损失使用以 2 为底的对数，则写作 $2^H$。PPL 越低表示模型平均给真实续词更高概率，“等效候选数”只是直觉，不是说模型真的只在固定数量 token 间选择。
\end{importantbox}

\begin{warningbox}{Perplexity 不能跨 tokenizer 直接比较}
不同 tokenizer 会改变序列长度和每步预测单位，同一文本的 token-level PPL 因而不可直接横比。PPL 还可能在通用文本上改善，却让安全、工具调用或长任务表现退化。slide 76 的历史趋势支持语言建模能力显著提高，但不证明 2023 年后的 Agent 能力可由 PPL 单独预测。生产系统需要任务成功、成本、安全和恢复能力等额外指标。
\end{warningbox}

\begin{knowledgebox}{读图：控制长度是在问因果问题}
若模型 A 比 B 输出更长，judge 又偏好更长回答，原始 win rate 混合了“内容更好”和“更长”两种效应。length-controlled AlpacaEval 用回归或匹配方法控制长度后再估模型效应。它减少一个已知混杂因素，却不能自动消除事实性、风格、位置和 judge 自我偏好等其他偏差。正确做法是把评估器当作需要持续校准的测量仪器。
\end{knowledgebox}

